% vim: tw=78 ai spell spelllang=en_gb sw=2
% syntax spell toplevel

\documentclass[10pt]{llncs}
\usepackage{color}
\usepackage{colortbl}
\usepackage{url}
\usepackage[T1]{fontenc}
\usepackage{graphicx}
\usepackage{listings}
\usepackage{hyperref}
\usepackage{multirow}
\usepackage{tikz}
\usepackage{xspace}
\pagestyle{plain}

\usetikzlibrary{arrows, positioning, decorations.pathreplacing}

\lstset{language=C,tabsize=2,showstringspaces=false,
	breaklines=true,breakatwhitespace=true,escapeinside={(*@}{@*)},
	captionpos=b,basicstyle=\sffamily\small}
%numbers=left,frame=l

\newcommand{\code}[2][]{\lstinline[columns=fullflexible, #1]$#2$}
\newcommand{\file}[2][]{\code[language={}, #1]{#2}}
\newcommand{\commit}[1][]{\texttt{#1}}

\newcommand{\xes}{\texttt{x86}\xspace}
\newcommand{\xesl}{\texttt{x86\_32}\xspace}
\newcommand{\xesh}{\texttt{x86\_64}\xspace}

\setcounter{tocdepth}{1}

\begin{document}

\title{Stable Patches and SUSE Trees}
\author{Jiri~Slaby}
\institute{Suse Labs \\
\email{jslaby@suse.cz}}

\maketitle

\begin{abstract}
  Here, we provide a short summary of the past and current habits of merging
  stable patches into the SUSE kernel trees. Starting with the sorted series
  section, the policy undergoes a change, so we will sum up what is happening
  in SLE15 now and what is our plan for future releases.
\end{abstract}

\section{Introduction} \label{sec:intro}
SUSE is maintaining kernels for their products in a git tree. Each product has
a dedicated branch called after it. Every such branch is based on some picked
upstream kernel version which is immutable over time. On the top of it, the
branches contain from tens (master and stable branches) to tens of thousands
(SLE, Leap) of patches we maintain in them. To be able to apply the patches,
we also store their list in a file called \file{series.conf}.

In master and stable branches, these patches are placed into
\file{series.conf} in an order depending on a person who did the backport. The
patches are divided into so called \emph{sections}. There are sections for
filesystems, x86, TTY, sound and alike. For tens of patches it is quite
bearable to maintain the patches in this form. We also used to use this schema
for SLE.  However, lately, we switched to a \emph{sorted section}.

\paragraph{Sorted section} is a section in \file{series.conf} where all the
patches are sorted according to their appearance in the upstream kernel. This
makes lives of developers easier by the order of magnitude. The main reason is
that there are no interdependencies between patches placed in different
sections. These dependencies created conflicts when applying patches and also
made the patches to differ from the upstream ones.

One exception were stable patches. In upstream, for each kernel release, some
later minor releases with non-intrusive fixes are released. So for example, we
have 4.12.1, 4.12.2, \dots\ for a released kernel 4.12. These stable releases
are composed of many patches and we add them in a single commit per file
manner (like any other patch). For all branches master, stable, SLE, Leap we
put these stable patches to the very front of the series file. Obviously, this
creates exactly the same interdependencies from the upstream patches as
non-sorted sections above.

With 4.12, only 14 stable version were released. After that we add stable
patches from the upstream stable list and other our tools like git-fixes. All
these are put to the sorted section, despite they have the same origin and
reasons as stable patches.

\section{Unwinding} \label{sec:unwinding}
\section{Unwinding in openSUSE \& SLE} \label{sec:suse_unw}
\section{Conclusion} \label{sec:concl}

\bibliographystyle{plain}
\bibliography{paper}

\end{document}
